\section{Experimental Section}

\subsection{Synthesis of Ti${_3}$C${_2}$T${_x}$ MXene}

Ti$_{3}$C$_{2}$T$_{x}$ MXene was prepared by selectively etching the Al layer from Ti$_{3}$AlC$_{2}$ MAX phase (purity $\geq$ 98\%, $\leq$ 400 mesh) using a LiF/HCl etchant. In a typical procedure, 3.2~g of LiF was dissolved in 40~mL of HCl (12~mol~L$^{-1}$) in a polytetrafluoroethylene (PTFE) beaker under magnetic stirring at room temperature for 30~min. Subsequently, 2~g of Ti$_{3}$AlC$_{2}$ powder was slowly added to the mixture, and the reaction was maintained at 40~$^{\circ}$C in a water bath with continuous stirring for 48~h. The resulting product was centrifuged at 5,000~rpm for 1~min, and the supernatant was discarded. The collected sediment was repeatedly washed with 1~mol~L$^{-1}$ HCl and deionized water by centrifugation (5,000~rpm, 1~min per cycle) until the pH of the supernatant approached neutral. The obtained MXene was then dispersed in 80~mL of deionized water and subjected to ice-bath ultrasonication for 30~min. After ultrasonication, the dispersion was centrifuged at 3,500~rpm for 15~min, and the dark supernatant — a high-concentration MXene colloidal suspension (20--25~mg~mL$^{-1}$) — was collected. The final MXene suspension was purged with N$_{2}$ and stored under refrigeration.

\subsection{Synthesis of Bi/MXene Composite}

The Bi/MXene composite was synthesized via a hydrothermal route. First, freeze-dried MXene aerogel (0.21~g) was added to 40~mL of anhydrous ethanol containing 2~mmol of BiCl$_{3}\cdot x$H$_{2}$O. The mixture was magnetically stirred for 3~h to ensure uniform dispersion, then transferred to a PTFE-lined stainless-steel autoclave and heated at 180~$^{\circ}$C for 16~h. During this process, Bi$^{3+}$ underwent hydrolysis to form Bi(OH)$_{3}$ nanoparticles that nucleated on the MXene surface, yielding Bi(OH)$_{3}$/MXene. After the reaction, the product was collected by centrifugation, washed several times with anhydrous ethanol and deionized water, and dried in a vacuum oven at 60~$^{\circ}$C. Finally, the dried powder was calcined at 350~$^{\circ}$C for 2~h under Ar atmosphere in a tube furnace to convert Bi(OH)$_{3}$ to metallic Bi, yielding the Bi/MXene composite.

For comparison, pristine Bi was prepared under identical conditions without the addition of MXene, using H$_{2}$/Ar atmosphere during calcination.

\subsection{Materials Characterization}

The morphology and microstructure of the samples were examined by field-emission scanning electron microscopy (SEM, JEOL JSM-6700F), equipped with an energy-dispersive X-ray spectrometer (EDS) for elemental mapping. High-resolution transmission electron microscopy (HRTEM, JEOL JEM-2010) was employed to analyze the atomic-scale interfacial structure. X-ray diffraction (XRD, Rigaku D/Max2500) with Cu K$\alpha$ radiation ($\lambda = 1.5406$~\AA) was used to determine the crystal phases. The chemical states of the elements were analyzed by X-ray photoelectron spectroscopy (XPS, Thermo Fisher Scientific ESCALAB 250Xi).

\subsection{Electrochemical Measurements}

Electrochemical measurements were performed using CR2032 coin-type half-cells assembled in an Ar-filled glove box (H$_{2}$O, O$_{2}$ < 0.1~ppm). The working electrode was prepared by casting a homogeneous slurry of the active material (Bi/MXene or pristine Bi), Super~P carbon black, sodium carboxymethyl cellulose (CMC), and styrene-butadiene rubber (SBR) in a weight ratio of 7:1:1:1 in deionized water onto a copper foil current collector. The coated foil was dried under vacuum at 60~$^{\circ}$C for 12~h. The mass loading of the active material was approximately 1~mg~cm$^{-2}$. Sodium metal foil served as both the counter and reference electrode, with a glass fiber separator (Whatman GF/D) and 1.0~mol~L$^{-1}$ NaPF$_{6}$ in dimethoxyethane (DME) as the electrolyte.

Galvanostatic charge--discharge (GCD) tests were conducted on a Neware battery testing system in the voltage range of 0.01--1.8~V (vs. Na$^{+}$/Na). Rate capability was evaluated at current densities from 0.1 to 20~A~g$^{-1}$. Long-term cycling stability was assessed at 2~A~g$^{-1}$ for 5,000 cycles. Cyclic voltammetry (CV) was performed on a Princeton Applied Research electrochemical workstation at scan rates of 0.2, 0.4, 0.6, 0.8, and 1.0~mV~s$^{-1}$ over the same voltage window. Electrochemical impedance spectroscopy (EIS) was recorded over a frequency range of 100~kHz to 0.01~Hz with an AC amplitude of 5~mV. The galvanostatic intermittent titration technique (GITT) was performed by applying a current pulse of 0.05~mA for 10~min followed by a relaxation period of 60~min.

For full-cell assembly, Na$_{3}$V$_{2}$(PO$_{4}$)$_{3}$ (NVP) powder was mixed with Super P and PVDF in a weight ratio of 8:1:1, ground, and dispersed in $N$-methyl-2-pyrrolidone. The homogeneous slurry was coated onto aluminum foil and dried under vacuum at 353~K overnight. Glass fiber (Whatman GF/D) and 1~M NaPF$_{6}$ in DME served as the separator and electrolyte, respectively. The cycling performance of NVP//Na half-cells was evaluated in the voltage range of 2.0--4.0~V. For NVP//Bi/MXene full cells, the Bi/MXene anode and NVP cathode were assembled in CR2032 coin cells at a mass ratio of 1:4. Prior to assembly, the Bi/MXene anode was pre-cycled for three charge--discharge cycles at 0.1~A~g$^{-1}$. The full-cell capacity was calculated based on the anode mass, and cycling performance was tested in the voltage window of 1.6--3.2~V.
